%PACOTES ABNT
\usepackage[brazil]{babel} % Pacote para usar portugues.
\usepackage[T1]{fontenc} % Pacote para copiar palavras com acentuação do pdf
\usepackage[utf8]{inputenc} % Pacote para poder usar acentuação no arquivo .tex

\usepackage{graphicx} % Pacote para inserir Figuras.
\usepackage[top=2.0cm,right=2.0cm,left=2.0cm,bottom=2.0cm]{geometry} % Pacote de Margens
\usepackage{enumerate} % Pacote para usar marcadores.
\usepackage[onehalfspacing]{setspace} % Pacote espaçamento 1,5 cm (Conforme Modelo Oficial CONICT rev05)
\usepackage[alf]{abntex2cite} % Configura o sistema de citações e referências para o estilo ABNT.

% PACOTES ESSENCIAIS
\usepackage{graphicx,xcolor,comment,enumerate,multirow,multicol,indentfirst}

% PACOTES DE MATEMÁTICA
\usepackage{amsmath,amsthm,amsfonts,amssymb,dsfont,mathtools,blindtext}

\usepackage{sectsty}
\usepackage{titlecaps}
\sectionfont{\normalsize\MakeUppercase}

\usepackage{fancyhdr} % Para adicionar o rodapé
\usepackage{microtype} % Ajuste fino para evitar estouros de linha
\usepackage{url}
\urlstyle{same}
